\documentclass{article}
\usepackage{amsmath, amssymb, amsthm}
\usepackage{geometry}
\geometry{a4paper, margin=1in}
\usepackage{fancyhdr}
\usepackage{hyperref}
\usepackage[english, russian]{babel}
\usepackage[utf8]{inputenc}
\usepackage{listings}
\usepackage{xcolor}
\usepackage{tikz}

\pagestyle{fancy}
\fancyhf{}
\fancyhead[L]{Algorithms}
\fancyhead[R]{\thepage}

\title{Algorithms 1 \\ Homework 10}
\author{Nikita Zagainov}
\date{\selectlanguage{english}\today}

\lstset{
    language=Python,
    basicstyle=\ttfamily\small,
    keywordstyle=\color{blue},
    stringstyle=\color{red},
    commentstyle=\color{gray},
    backgroundcolor=\color{lightgray!20},
    frame=single,
    rulecolor=\color{black},
    numbers=left,
    numberstyle=\tiny\color{gray},
    stepnumber=1,
    numbersep=10pt,
    showspaces=false,
    showstringspaces=false,
    breaklines=true,
    breakatwhitespace=true,
    tabsize=4,
    captionpos=b
}

\begin{document}

\maketitle

\begin{center}
	\section*{Доказательства утверждений в графах.
	  Кратчайшие пути. Разные задачи.}
\end{center}
\begin{enumerate}
	\item Для решения этой задачи подходит алгоритм
	      Дейкстры. Добавление константной величины к длине
	      пути в каждой вершине не влияет на результат. Докажем
	      это: при переходе в некоторую вершину, чтобы учесть
	      константную величину, добавляемую к длине пути,
	      можно добавить её к весу ребра, ведущего в эту вершину.
	      А так как эта константа одинакова для всех вершин,
	      то порядок обхода не изменится. Тогда алгоритм Дейкстры
	      даст корректный результат.

	      Алгоритм:
	      \begin{lstlisting}
import heapq


def dijkstra(
    graph: dict[object, tuple[float, object]], start: object, end: object
) -> tuple[float, list[object]]:
    queue = [(0, [start])]
    visited = set()
    while queue:
        cost, path = heapq.heappop(queue)
        cur_node = path[-1]
        visited.add(cur_node)
        if cur_node == end:
            return cost, path

        for weight, neighbor in graph[cur_node]:
            if neighbor not in visited:
                heapq.heappush(
                    queue,
                    (
                        cost + weight + 3,
                        path + [neighbor],
                    ),
                )
    \end{lstlisting}

	      Сложность - $O((|V| + |E|) \log |V|) =
		      O((|V| + |V|k) \log |V|)$, так как
	      количество коридоров по условию не больше $\frac{|V|k}{2}$

	\item
	      \begin{enumerate}
		      \item
		            Предлагаемый алгоритм - алгоритм поиска компонент
		            сильной связности в данном графе:
		            \begin{lstlisting}
def inverse_graph(graph: dict[object, list[object]]) -> dict[object, list[object]]:
	reversed_graph = {}

	for node in graph:
		for neighbor in graph[node]:
			if neighbor not in reversed_graph:
				reversed_graph[neighbor] = list()
			reversed_graph[neighbor].append(node)

	return reversed_graph


def _dfs(
	graph: dict[object, list[object]],
	node: object,
	visited: set[object],
	start: list[object],
	finish: list[object],
):
	visited.add(node)
	start.append(node)
	finish.append(None)
	for neighbor in graph.get(node, []):
		if neighbor not in visited:
			_dfs(graph, neighbor, visited, start, finish)
	start.append(None)
	finish.append(node)


def dfs(
	graph: dict[object, list[object]],
) -> tuple[list[object | None], list[object | None]]:
	visited = set()
	start = list()
	finish = list()

	for node in graph:
		if node in visited:
			continue
		_dfs(graph, node, visited, start, finish)

	return start, finish


def graph_condensates(graph: dict[object, list[object]]) -> list[list[object]]:
	_, _finish = dfs(inverse_graph(graph))
	order = [node for node in reversed(_finish) if node is not None]

	visited = set()
	condensates = list()

	start = list()
	finish = list()

	for node in order:
		if node in visited:
			continue
		_dfs(graph, node, visited, start, finish)
		condensates.append([n for n in finish if n is not None])
		start.clear()
		finish.clear()

	return condensates		
	\end{lstlisting}

		            Сложность алгоритма ---
		            $O(\left| V \right| + \left| E \right|)$,
		            так как используется инвертирование графа за
		            $O(\left| V \right| + \left| E \right|)$ и
		            поиск в глубину на инвертированном графе за
		            $O(\left| V \right| + \left| E \right|)$.
		            Этот алгоритм корректен, так как он находит
		            компоненты сильной связности в графе по построению.

		      \item Для решения этой задачи, рассмотрим граф из
		            компонент сильной связности, построенный в предыдущем пункте.
		            этот граф не имеет циклов. Обозначим вершины этого графа
		            как истоки, если к ним не ведёт ни одно ребро из других
		            и как стоки, если из них не ведёт ни одно ребро в другие.
		            Очевидно, что минимальное количество
		            новых рёбер не больше максимума из количества стоков и истоков.
		            Докажем, что количество новых рёбер может равняться этому числу.
		            Полученный после применения поиска КСС граф является множеством
		            подграфов, в которых в сумме $n$ стоков и $m$ стоков.
		            Пока есть несвязные компоненты, проводим рёбра из стоков
		            произвольных подграфов в истоки тех подграфов, если подграфов
		            всего $k \leq \min(n, m)$, то в результате этого действия
		            останется связный граф с $n-k$ стоков и $m-k$ истоков.
		            Потребуется ещё $\max(n, m) - k$ рёбер, чтобы избавится от
		            лишних стоков и истоков.

		            Алгоритм:
		            \begin{lstlisting}
def get_sources_sinks(
	graph: dict[object, set[object]],
) -> tuple[set[object], set[object]]:
	all_nodes = set(graph.keys())
	all_neighbors = set(
		neighbor for neighbors in graph.values() for neighbor in neighbors
	)
	sources = all_nodes - all_neighbors
	sinks = all_neighbors - all_nodes

	return sources, sinks


def add_edges_to_connected(
	graph: dict[object, list[object]],
) -> dict[object, list[object]]:
	condensates = graph_condensates(graph)
	if len(condensates) == 1:
		return graph

	node_to_condensate = dict()
	condensate_to_node = dict()
	for idx, condensate in enumerate(condensates):
		for node in condensate:
			node_to_condensate[node] = idx

		condensate_to_node[idx] = node

	condensed_graph = dict()
	for source, neighbors in graph.items():
		source_condensate = node_to_condensate[source]
		for neighbor in neighbors:
			destination_condensate = node_to_condensate[neighbor]
			if (
				source_condensate != destination_condensate
				and destination_condensate
				not in condensed_graph.get(source_condensate, set())
			):
				if source_condensate not in condensed_graph:
					condensed_graph[source_condensate] = set()
				condensed_graph[source_condensate].add(
					destination_condensate
				)

	sources, sinks = get_sources_sinks(condensed_graph)

	while len(sources) > 0 and len(sinks) > 0:
		source = next(iter(sources))
		visited = set()
		_dfs(condensed_graph, source, visited, [], [])
		possible_sinks = sinks.difference(visited)

		if len(possible_sinks) == 0:
			sink = next(iter(sinks))
		else:
			sink = next(iter(possible_sinks))

		sources.remove(source)
		sinks.remove(sink)

		if condensate_to_node[sink] not in graph:
			graph[condensate_to_node[sink]] = []
		graph[condensate_to_node[sink]].append(
			condensate_to_node[source]
		)

	while len(sinks) > 0:
		sink = next(iter(sinks))
		sinks.remove(sink)
		if condensate_to_node[sink] not in graph:
			graph[condensate_to_node[sink]] = []
		graph[condensate_to_node[sink]].append(
			condensate_to_node[source]
		)

	while len(sources) > 0:
		source = next(iter(sources))
		sources.remove(source)
		if condensate_to_node[sink] not in graph:
			graph[condensate_to_node[sink]] = []
		graph[condensate_to_node[sink]].append(
			condensate_to_node[source]
		)

	return graph	
				\end{lstlisting}

		            Сложность алгоритма - $O(\left| V \right| + \left| E \right|)$.
		            Первым делом он за $O(\left| V \right| + \left| E \right|)$
		            находит конденсаты, далее за
		            $O(\left| V \right| + \left| E \right|)$
		            он строит граф конденсатов, а затем за
		            $O(\left| V \right| + \left| E \right|)$
		            с помощью нескольких вызовов поиска в глубину,
		            (суммарный размер обхода ограничен суммарным
		            размером подграфов), добавляет новые рёбра.
		            Если же остались неспаренные стоки или истоки,
		            он их добавляет к циклу за $O(\left| V \right|)$.
	      \end{enumerate}

	\item Для решения этой задачи можно обойти граф, раскрашивая
	      вершины в два цвета. Если обнаруживается, что две соседние
	      вершины имеют одинаковый цвет, то граф имеет цикл нечётной
	      длины. Если же такого не происходит, то граф не имеет циклов.
	      Докажем это утверждение: если граф не имеет циклов, то этот
	      граф --- дерево, и его можно раскрасить в два цвета. Граф,
	      содержащий циклы, может быть получен путём добавления
	      рёбер к дереву. Соединив две вершины дерева, мы получим
	      цикл нечётной длины, если эти вершины имеют одинаковый цвет
	      и цикл чётной длины, если вершины имеют разные цвета.
	      Таким образом утверждение верно по построению.

	      Алгоритм:
	      \begin{lstlisting}
def has_neg_cycle(graph: dict[object, list[object]]) -> bool:
    colors = dict()
    stack = list(graph.keys())

    while stack:
        node = stack.pop()
        if node not in colors:
            colors[node] = 0

        for neighbor in graph.get(node, []):
            if neighbor not in colors:
                colors[neighbor] = 1 - colors[node]
                stack.append(neighbor)
            elif colors[neighbor] == colors[node]:
                return True

    return False
        \end{lstlisting}
	      Сложность алгоритма --- $O(|V| + |E|)$.

	\item Алгоритм обхода в глубину позволяет найти кратчайшие
	      пути в графе с единичными весами рёбер, так как в таком
	      графе порядок обхода вершин будет совпадать с обходом
	      алгоритма Дейкстры. Теперь докажем, что в графе состояний
	      кодового замка для каждой вершины существует другая вершина,
	      расстояние до которой будет равно диаметру графа. Пусть
	      диаметр совпадает с кратчайшем путём между вершинами
	      $abcd$ и $efgh$. Кратчайшим путём между вершинами $abcd$ и
	      $efgh$ будет равен $|a - e| + |b - f| + |c - g| + |d - h|$.
	      тогда повернув кодовый замок на одно деление в произвольной
	      позиции в обоих замках, мы получим две новые вершины, расстояние между
	      которыми будет равно
	      $|(a + 1) \mod 10 - (e + 1) \mod 10| +
		      |b - f| + |c - g| + |d - h| =
		      |a - e| + |b - f| + |c - g| + |d - h|$. Данной операцией
	      мы можем получить любую вершину. Тогда диаметр графа
	      можно найти обходом в ширину из любой вершины.

	      Алгоритм:
	      \begin{lstlisting}
from collections import deque


def switch_lock(lock: str, index: int, direction: int) -> str:
	lock_list = list(lock)
	lock_list[index] = str((int(lock_list[index]) + direction) % 10)
	return "".join(lock_list)


def generate_lock_graph() -> dict[str, list[str]]:
	graph = {}
	for i in range(10000):
		lock_value = f"{i:04d}"
		graph[lock_value] = []
		for j in range(4):
			new_value_1 = switch_lock(lock_value, j, 1)
			new_value_2 = switch_lock(lock_value, j, -1)
			graph[lock_value].append(new_value_1)
			graph[lock_value].append(new_value_2)
	return graph


def get_lock_graph_diameter() -> int:
	graph = generate_lock_graph()
	visited = set()
	queue = deque([(0, "0000")])

	max_distance = 0
	while queue:
		distance, lock = queue.popleft()
		max_distance = max(max_distance, distance)
		visited.add(lock)

		for next_lock in graph[lock]:
			if next_lock in visited:
				continue
			queue.append((distance + 1, next_lock))

	return max_distance
	\end{lstlisting}
	      Сложность алгоритма --- $O(|V| + |E|)$.

	\item Для нахождения кратчайшего пути в невзвешенном графе
	      можно воспользоваться алгоритмом $BFS$. Вызываем его для
	      каждого владельца таксы, чтобы найти ближайшего хозяина таксы
	      для каждого их них.

	      Алгоритм:
	      \begin{lstlisting}
from collections import deque

def find_closest_dog_masters(
    graph: dict[object, list[object]], dog_owners: set[object]
) -> dict[object, list[object]]:
    if len(dog_owners) == 1:
        return {}
    
    result = {}
    for dog_owner in dog_owners:
        queue = deque([[dog_owner]])
        visited = set()
        while queue:
            chain = queue.popleft()
            cur_node = chain[-1]
            visited.add(cur_node)
            if cur_node != dog_owner and cur_node in dog_owners:
                result[cur_node] = chain
                break

            for neighbor in graph[cur_node]:
                if neighbor in visited:
                    continue

                queue.append(chain + [neighbor])

    return result
		  \end{lstlisting}
		  Сложность --- $O(m(|V| + |E|))$.
\end{enumerate}

\end{document}


